\begin{lemma}[The concave radial profile $f(r)=r-\tfrac18r^{3/2}$]
Put
\[
  f(r)=r-\tfrac18 r\sqrt r\qquad(r\ge0),
\]
so that $f(r)=r-\tfrac18r^{3/2}$, put $f_1(r)=1-\tfrac3{16}\sqrt r$ and put
$\nu=1-\tfrac18\sqrt{3/2}$. Then
\begin{enumerate}
\item $\nu>0$;
\item $f$ is of class $C^2$ at every $r>0$;
\item for every $r>0$, $f$ is differentiable at $r$ with $f'(r)=f_1(r)$;
\item for every $r>0$, $f_1$ is differentiable at $r$ with
  $f_1'(r)=-\dfrac3{32\sqrt r}$ (so $f''(r)=-3/(32\sqrt r)$);
\item for every $r$ with $0<r\le\tfrac32$,
  \[
    \nu r\le f(r)\le r,\qquad 0<f_1(r)\le1 .
  \]
\end{enumerate}
\end{lemma}
\begin{proof}
Throughout, $\sqrt{\cdot}$ denotes the nonnegative square root, and for $r>0$ we
have $\sqrt r>0$ and $(\sqrt r)^2=r$.

\emph{A numerical bound used twice.} Since $3/2<4$ and the square root is
strictly increasing on $[0,\infty)$, $\sqrt{3/2}<\sqrt4=2$. In particular
\begin{equation}\label{eq:crp-num}
  0\le\sqrt{3/2}<2 .
\end{equation}

\emph{(1).} By \eqref{eq:crp-num}, $\tfrac18\sqrt{3/2}<\tfrac28=\tfrac14<1$, so
$\nu=1-\tfrac18\sqrt{3/2}>0$.

\emph{(2).} Fix $r>0$. The square root is $C^\infty$ on $(0,\infty)$, hence
$C^2$ at $r$; the identity $s\mapsto s$ is $C^\infty$; a product of two
functions that are $C^2$ at a point is $C^2$ there, and so are a constant
multiple of such a function and a difference of two of them. Hence
$f(s)=s-\tfrac18 s\sqrt s$ is $C^2$ at $r$.

\emph{(3).} Fix $r>0$. The square root is differentiable at $r$ with derivative
$1/(2\sqrt r)$. Hence, by the product rule applied to $s\mapsto s$ and
$s\mapsto\sqrt s$, the function $s\mapsto s\sqrt s$ is differentiable at $r$
with derivative
\[
  1\cdot\sqrt r+r\cdot\frac1{2\sqrt r}
   =\sqrt r+\frac{(\sqrt r)^2}{2\sqrt r}
   =\sqrt r+\frac{\sqrt r}{2}
   =\frac32\sqrt r ,
\]
where $r=(\sqrt r)^2$ was used in the middle step. Dividing by the constant $8$
and subtracting from the identity $s\mapsto s$, whose derivative is $1$, shows
that $f$ is differentiable at $r$ with
\[
  f'(r)=1-\frac1 8\cdot\frac32\sqrt r=1-\frac3{16}\sqrt r=f_1(r).
\]

\emph{(4).} Fix $r>0$. Again the square root is differentiable at $r$ with
derivative $1/(2\sqrt r)$; multiplying by the constant $3$, dividing by the
constant $16$ and subtracting from the constant $1$ shows that $f_1$ is
differentiable at $r$ with
\[
  f_1'(r)=-\frac{3}{16}\cdot\frac1{2\sqrt r}=-\frac3{32\sqrt r} .
\]

\emph{(5).} Fix $r$ with $0<r\le\tfrac32$. Monotonicity of the square root gives
$0<\sqrt r\le\sqrt{3/2}$.

First, $f(r)=r\bigl(1-\tfrac18\sqrt r\bigr)$. Since $\sqrt r\le\sqrt{3/2}$ we
have $1-\tfrac18\sqrt r\ge1-\tfrac18\sqrt{3/2}=\nu$, and $r>0$, so
$f(r)\ge\nu r$. Since $\sqrt r>0$ we have $1-\tfrac18\sqrt r\le1$, so
$f(r)\le r$.

Second, by \eqref{eq:crp-num}, $\sqrt r\le\sqrt{3/2}<2$, hence
$\tfrac3{16}\sqrt r<\tfrac6{16}<1$ and so $f_1(r)=1-\tfrac3{16}\sqrt r>0$;
and $\sqrt r>0$ gives $f_1(r)\le1$.
\end{proof}
